% TOP_PROBLEM: {"id": 30006882, "title": "Strong Hyperplane Conjecture for Isotropic Constants", "category_id": 6, "category": {"id": 6, "name": "geometry", "display_name": "Geometry", "description": "Euclidean and non-Euclidean geometry, geometric structures.", "slug": "geometry", "order_index": 6, "created_at": "2026-07-31T15:26:25.671Z"}, "status": "open"}
% FIRST_POSTED: 2026-09-27
% !TeX program = lualatex
% ATTEMPT_STATUS: unresolved
% Model: GPT-6 Astra Ultra; continuation dated 27 September 2026.
\documentclass[11pt]{article}
\usepackage[margin=25mm]{geometry}
\usepackage{fontspec}
\setmainfont{Latin Modern Roman}
\usepackage{amsmath,amssymb,mathtools}
\usepackage{unicode-math}
\setmathfont{Latin Modern Math}
\usepackage{xurl,hyperref}
\hypersetup{colorlinks=true,urlcolor=blue}
\setcounter{secnumdepth}{0}
\setlength{\parindent}{0pt}
\setlength{\parskip}{5pt}
\setlength{\emergencystretch}{3em}
\begin{document}
\section*{\texorpdfstring{Strong Hyperplane Conjecture for Isotropic Constants}{Strong Hyperplane Conjecture for Isotropic Constants}}\label{30006882}
\subsection{Definitions and mathematical statement}
For a convex body $K\subset{\mathbb{R}}^n$, let $V=|K|$, $b=V^{-1}\int_Kx{\,\mathrm{d}} x$, and
\[
 \Sigma_K=\frac1V\int_K(x-b)(x-b)^{\mathsf T}{\,\mathrm{d}} x,
 \qquad L_K=\left(\frac{\det\Sigma_K}{V^2}\right)^{1/(2n)}.
\]
The covariance is positive definite because $K$ has nonempty interior. This definition is invariant under invertible affine changes of coordinates. The sharp conjecture is
\[
 L_K\le L_{\Delta_n}
 =\frac{(n!)^{1/n}}{(n+1)^{(n+1)/(2n)}\sqrt{n+2}},
\]
where $\Delta_n$ is any nondegenerate simplex. It applies in each dimension $n\ge1$, without symmetry or smoothness assumptions. It seeks the simplex's precise dimension-dependent bound, not merely some universal upper bound or a uniqueness statement for equality.
\par\smallskip\textit{Formulation and concept references:} \href{https://arxiv.org/pdf/2407.01353v2}{[S1]}.

\subsection{Short English statement}
In every dimension, does a simplex maximize the affine-invariant isotropic constant among all convex bodies?
\subsection{Research attempt}
\paragraph{Which hyperplane conjecture is being tested.}
The statement fixes a dimension-dependent extremal value. It must
not be confused with the existence of some dimension-free upper
bound. The inspected revision [S1], dated 27 May 2025, explicitly
distinguishes the latter result from the strong simplex conjecture.
Its Theorem 1.2 concerns polytopal local maximizers with a simplicial
vertex; it does not say every maximizer has such a vertex. We retain
that distinction. The arguments below compute several model families,
derive a first-variation obstruction, and explain a concrete failure
of a proposed reduction by approximation.

Write $F(K)=L_K^{2n}=\det\Sigma_K/|K|^2$. If $T(x)=Ax+c$
with $A$ invertible, then
\[
 |T(K)|=|\det A||K|,\quad b_{T(K)}=Ab_K+c,
 \quad\Sigma_{T(K)}=A\Sigma_KA^{\mathsf T}.
\]
Consequently $F(T(K))=F(K)$. In particular, translating by the
barycenter and applying $\Sigma_K^{-1/2}$ puts the body in the
normalization $b=0$, $\Sigma=I$. In this normalization
\[
 L_K=|K|^{-1/n},\qquad \int_K|x|^2\,dx=n|K|.
 \tag{358.1}
\]
This is covariance-one normalization, rather than the alternative
volume-one normalization sometimes also called isotropic position.

\paragraph{The simplex value, including its determinant.}
Take the coordinate simplex
$S=\{x_i\geq0:\sum_{i=1}^nx_i\leq1\}$. Repeated integration,
or induction using the last coordinate, gives
\[
 \int_Sx_1^{a_1}\cdots x_n^{a_n}\,dx
  =\frac{\prod_{i=1}^na_i!}{(n+\sum_i a_i)!}
 \qquad(a_i\hbox{ nonnegative integers}).
 \tag{358.2}
\]
For the induction, integrate the first $n-1$ coordinates over a
simplex of scale $1-x_n$ and use
$\int_0^1t^a(1-t)^b\,dt=a!b!/(a+b+1)!$, which itself follows
by repeated integration by parts. Thus (358.2) does not assume a
probability formula for simplex coordinates.

It follows that $|S|=1/n!$, $b_i=1/(n+1)$ and
\[
 \mathbb E X_i^2=\frac2{(n+1)(n+2)},\qquad
 \mathbb E X_iX_j=\frac1{(n+1)(n+2)}\ (i\ne j).
\]
Subtracting the means gives
\[
 \Sigma_S=\frac{(n+1)I-J}{(n+1)^2(n+2)},
 \qquad
 \det\Sigma_S=\frac1{(n+1)^{n+1}(n+2)^n},
 \tag{358.3}
\]
where $J$ is the all-ones matrix. To check the determinant, the
numerator has eigenvalue one in the all-ones direction and eigenvalue
$n+1$ on its orthogonal complement. Combining (358.3) with the volume
recovers exactly the printed simplex constant. All nondegenerate
simplexes have the same value by affine invariance.

For $n=1$ every convex body with interior is an interval. Its length
$\ell$ and variance $\ell^2/12$ give $L=1/\sqrt{12}$, so the
selected conjecture is proved completely in dimension one, including
all translations and lengths.

\paragraph{Product bodies and two comparison families.}
For full-dimensional $K\subset\mathbb R^a$ and
$H\subset\mathbb R^b$, the uniform random point of $K\times H$
has independent components. Its covariance is block diagonal and
its volume is the product. Therefore
\[
 L_{K\times H}^{a+b}=L_K^aL_H^b.
 \tag{358.4}
\]
In particular every parallelepiped has $L=1/\sqrt{12}$. Product
structure averages logarithms of isotropic constants; it cannot
produce a value larger than both factors. But this identity does
not reduce arbitrary bodies to products: most bodies have no such
affine decomposition.

For the Euclidean unit ball $B_2^n$, write its volume as $\kappa_n$.
Rotational invariance gives a scalar covariance matrix, and radial
integration gives $\mathbb E|X|^2=n/(n+2)$. Hence
\[
 L_{B_2^n}=\frac1{\sqrt{n+2}\,\kappa_n^{1/n}}.
 \tag{358.5}
\]
For the cross-polytope $B_1^n=\{\sum_i|x_i|\leq1\}$, the
$2^n$ orthants are reflected coordinate simplexes. Reflection makes
all means and mixed second moments vanish, while (358.2) gives
$\mathbb E X_i^2=2/((n+1)(n+2))$ and $|B_1^n|=2^n/n!$.
Thus
\[
 L_{B_1^n}=\frac{(n!)^{1/n}}{\sqrt{2(n+1)(n+2)}}.
 \tag{358.6}
\]
The ratio to the simplex value is
\[
 \frac{L_{B_1^n}}{L_S}
   =\frac{(n+1)^{1/(2n)}}{\sqrt2}\leq1,
 \tag{358.7}
\]
because $n+1\leq2^n$ for $n\geq1$, by elementary induction.
This proves the desired comparison for every affine cross-polytope.
It is a full family calculation, not a numerical comparison in low
dimensions. In dimension two the cross-polytope is a parallelogram,
consistent with (358.4).

\paragraph{An extremal statement in the opposite direction.}
There is a direct proof that ellipsoids minimize this functional. It
is worth spelling out because a moment-minimization argument can
otherwise be accidentally used with the wrong sign.

Normalize an arbitrary $K$ as in (358.1) and set $V=|K|$. Let $D$
be the centered ball of volume $V$, with radius
$r=(V/\kappa_n)^{1/n}$. Since $|K\setminus D|=|D\setminus K|$,
\[
 \int_K|x|^2-\int_D|x|^2
 =\int_{K\setminus D}(|x|^2-r^2)
   +\int_{D\setminus K}(r^2-|x|^2)\geq0.
\]
Using $\int_K|x|^2=nV$ and
$\int_D|x|^2=nVr^2/(n+2)$ gives $r^2\leq n+2$. Therefore
\[
 V\leq\kappa_n(n+2)^{n/2},\qquad
 L_K\geq\frac1{\sqrt{n+2}\kappa_n^{1/n}}.
 \tag{358.8}
\]
Equality forces $K=D$ up to measure zero, and convex bodies with
interior agreeing almost everywhere agree as closed sets. Undoing
the affine normalization proves the ellipsoid equality statement.
The inequality points downward for $L_K$. Rearrangement supplies
the minimum, whereas the question asks for the maximum.

\paragraph{First variation under a genuine smooth deformation.}
Suppose temporarily that $K$ has a smooth boundary with strictly
positive principal curvatures. Normalize $b=0$ and $\Sigma=I$.
Let $K_t$ be a smooth normal deformation with scalar normal velocity
$h$ on the boundary. Denote volume, first moment and second moment by
$V(t)$, $M(t)=\int_{K_t}x\,dx$, $Q(t)=\int_{K_t}xx^{\mathsf T}\,dx$.
Differentiating integrals over a moving boundary gives
\[
 V'=\int_{\partial K}h\,dS,\quad
 M'=\int_{\partial K}xh\,dS,\quad
 Q'=\int_{\partial K}xx^{\mathsf T}h\,dS.
\]
These formulas follow by integrating the thickness $th+o(t)$ of the
boundary layer in local normal coordinates. Since $M(0)=0$, the
derivative of $bb^{\mathsf T}$ is zero at $t=0$. As $Q(0)=VI$,
\[
 \Sigma'=V^{-1}Q'-(V'/V)I.
\]
Jacobi's formula for the determinant now gives
\[
 \left.\frac{d}{dt}\right|_{0}\log F(K_t)
 =\frac1V\int_{\partial K}(|x|^2-n-2)h(x)\,dS(x).
 \tag{358.9}
\]
The centering correction has been included rather than silently held
fixed; it drops out only because its first variation vanishes at
zero mean. The coefficient $n+2$ comes from both covariance
normalization and the squared volume in $F$.

Every smooth $h$ is available in both signs for sufficiently small
$|t|$ under the positive-curvature assumption. One way to see this
is to perturb the support function on the sphere: its curvature
matrix is positive definite, and positivity survives a sufficiently
small $C^2$ perturbation. The Gauss map identifies the support-function
velocity with the desired normal velocity.

If such a body were a local maximizer, (358.9) would vanish for every
$h$, forcing $|x|^2=n+2$ on its boundary. Since a centered convex body
has zero in its interior, every ray meets its boundary once; the
constant radius then forces it to be the ball of radius $\sqrt{n+2}$.
For $n\geq2$ that ball is not a local maximum: (358.8) is strict for
every nearby nonellipsoid. Such nearby strictly convex bodies exist,
for example through a small smooth support perturbation outside the
finite-dimensional family of ellipsoids. We have therefore proved:
\[
 \text{A smooth positively curved body in dimension }n\geq2
 \text{ cannot locally maximize }L_K.
 \tag{358.10}
\]
This does not exclude the conjectured simplex; its boundary is outside
the regularity and two-sided deformation assumptions of the argument.

\paragraph{Why this variation does not classify nonsmooth maximizers.}
At a flat facet, an arbitrary normal displacement need not preserve
convexity for both signs. At an edge or a vertex the admissible
variations are more constrained still. Applying (358.9) as though all
boundary velocities were freely available would incorrectly rule out
polytopes, including the simplex. One must characterize a cone of
admissible variations and derive one-sided inequalities on that cone.
The source's polytopal theorems use additional geometric structure;
their hypotheses cannot be deleted by the smooth calculation above.

An approximation argument does not repair this. If convex bodies
$K_j$ converge to a full-dimensional $K$ in Hausdorff distance, their
volumes and first two moments converge. Indeed they lie in a common
bounded set and their indicators converge almost everywhere off
$\partial K$, which has Lebesgue measure zero; dominated convergence
then applies to the bounded coordinate monomials. Thus $L_{K_j}\to L_K$.
This allows upper bounds valid for \emph{every} approximating polytope
to pass to the limit. A theorem only about polytopes that locally
maximize $L$ is different: an arbitrary polytope in an approximating
sequence need not be a local maximizer. Smoothing has the same
problem in reverse, as (358.10) illustrates.

Likewise, triangulating a polytope does not express its covariance
as a volume-weighted average of the simplex covariance determinants.
For a disjoint union with component masses $w_i$, means $b_i$ and
covariances $\Sigma_i$, direct expansion gives
\[
 \Sigma=\sum_iw_i\Sigma_i+
          \sum_iw_i(b_i-b)(b_i-b)^{\mathsf T}.
 \tag{358.11}
\]
The positive between-component term is essential. The determinant
is nonlinear, and omitting this term would already give a wrong
answer for two adjacent intervals. Thus known simplex moments do
not prove the upper bound for a triangulated body.

\paragraph{Diagnostics and outcome.}
The accompanying exact checks evaluate the simplex moment matrix
and determinant, the cross-polytope ratio, and the product identity
in low dimensions. They test normalization and algebra only; no
finite list of shapes certifies a supremum over all convex bodies.
The actual proved results are the simplex value, the complete
one-dimensional case, the affine cross-polytope comparison, the
ellipsoid minimum, and the smooth positive-curvature obstruction.

The first unresolved step toward the selected sharp upper bound is
controlling possible nonsmooth maximizers without adding a simplicial
vertex, product decomposition, or unrestricted deformation assumption.
The universal slicing bound does not supply this sharp comparison.
The conjecture remains unresolved by this attempt.
\subsection{Sources}
\begin{itemize}
\item[S1]Isotropic constants and regular polytopes; arXiv2407.01353v2,27May2025.
\url{https://arxiv.org/pdf/2407.01353v2}.
\textit{Formulation source.}
\end{itemize}
\end{document}
